Program ditulis dalam Python dan disimpan di folder \texttt{src/}. Model
\textit{from scratch} hanya memakai NumPy~\cite{numpy2026docs}. Scikit-learn dipakai untuk pembagian
data, perhitungan metrik, dan model pembanding~\cite{sklearn2026docs}. Alur
kerja mengikuti tahapan CRISP-DM dan dijalankan dari \texttt{src/main.ipynb}.
Potongan kode pada bab ini diambil langsung dari \textit{repository}.

\section{Struktur Program}

\begin{table}[H]
	\centering
	\caption{Modul program}
	\label{tab:modul}
	\begin{tabular}{l p{7.6cm}}
		\toprule
		\textbf{Modul} & \textbf{Isi} \\
		\midrule
		\texttt{config.py} & Konstanta global (\textit{seed}, \textit{path}, nama kolom, rasio pembagian data, jumlah \textit{fold}) \\
		\texttt{eda.py} & Fungsi analisis eksploratif: tipe data, nilai hilang, duplikat, histogram, \textit{boxplot}, pencilan IQR, korelasi \\
		\texttt{preprocessing.py} & Pembagian data, \texttt{Preprocessor}, rekayasa fitur, PCA, dan \texttt{Pipeline} (seleksi fitur, PCA, model) \\
		\texttt{validation.py} & \textit{Hold-out}, \textit{k-fold}, \textit{tuning}, \textit{confusion matrix}, analisis kesalahan, \textit{resampling}, dan pembuatan berkas prediksi \\
		\texttt{models/base.py} & Antarmuka \texttt{BaseModel}: \texttt{fit}, \texttt{predict}, \texttt{set\_params}, \texttt{save}, \texttt{load} \\
		\texttt{models/knn.py} & K-Nearest Neighbors \\
		\texttt{models/logistic\_regression.py} & Logistic Regression \\
		\texttt{models/svm.py} & Support Vector Machine (SMO) \\
		\texttt{models/dtl.py} & Decision Tree Learning \\
		\texttt{models/bayes\_network.py} & Bayesian Network \\
		\texttt{main.ipynb} & \textit{Notebook} utama yang menjalankan seluruh tahapan CRISP-DM \\
		\bottomrule
	\end{tabular}
\end{table}

\section{Antarmuka Model}

Semua model \textit{from scratch} merupakan turunan \texttt{BaseModel}.
\texttt{fit(X, y)} menerima \textit{array} NumPy dan mengembalikan \texttt{self},
sedangkan \texttt{predict(X)} mengembalikan label 0 atau 1. Antarmuka ini sama
dengan antarmuka scikit-learn. Akibatnya, fungsi validasi, \textit{tuning}, dan
\texttt{Pipeline} dapat dipakai, baik untuk model \textit{from scratch} maupun
model pembanding, tanpa perubahan. Model yang sudah dilatih disimpan dan dimuat
kembali dengan \texttt{pickle}.

\lstinputlisting[caption={Antarmuka model (\texttt{src/models/base.py})},linerange={1-3,5-24}]{../src/models/base.py}

\section{K-Nearest Neighbors}
\label{sec:knn}

\todo{Implementasi \texttt{KNN.fit} dan \texttt{KNN.predict} di
\texttt{src/models/knn.py} belum selesai. Parameter yang sudah disiapkan:
\texttt{k}, \texttt{metric} (\texttt{euclidean}, \texttt{manhattan},
\texttt{minkowski}), dan \texttt{p}. Jelaskan cara kerja, fungsi/kelas, dan
sertakan potongan kode setelah implementasi selesai.}

\section{Logistic Regression}
\label{sec:lr}

\todo{Implementasi \texttt{LogisticRegression.fit} dan \texttt{predict} di
\texttt{src/models/logistic\_regression.py} belum selesai. Parameter yang sudah
disiapkan: \texttt{lr}, \texttt{n\_iter}, \texttt{reg} (L1/L2), dan \texttt{lam}.
Jelaskan fungsi sigmoid, \textit{log-loss}, \textit{gradient descent}, dan
regularisasi beserta potongan kode.}

\section{Support Vector Machine}
\label{sec:svm}

\subsection{Formulasi}

SVM dengan \textit{soft margin} mencari \textit{hyperplane} pemisah dengan margin
terbesar dan mengizinkan sebagian sampel melanggar margin.
Implementasi ini menyelesaikan bentuk dualnya. Untuk label $y_i\in\{-1,+1\}$ dan
fungsi kernel $K$, masalah dualnya adalah
\begin{equation}
	\max_{\alpha}\ \sum_{i}\alpha_i-\frac{1}{2}\sum_{i}\sum_{j}\alpha_i\alpha_j y_i y_j K(x_i,x_j)
	\quad\text{dengan}\quad 0\le\alpha_i\le C,\ \ \sum_i\alpha_i y_i=0 .
\end{equation}
Parameter $C$ menentukan besarnya hukuman untuk pelanggaran margin. Setelah
$\alpha$ diperoleh, prediksi dihitung dari \textit{support vector}, yaitu sampel
dengan $\alpha_i>0$:
\begin{equation}
	f(x)=\sum_{i\in SV}\alpha_i y_i K(x_i,x)+b,\qquad \hat{y}=\begin{cases}1 & f(x)\ge 0\\ 0 & f(x)<0\end{cases}
\end{equation}
Implementasi mendukung empat kernel (Tabel~\ref{tab:kernel}). Nama \texttt{rbf}
dan \texttt{gaussian} merujuk ke kernel yang sama.

\begin{table}[H]
	\centering
	\caption{Kernel SVM yang didukung}
	\label{tab:kernel}
	\begin{tabular}{l l}
		\toprule
		\textbf{Kernel} & $K(a,b)$ \\
		\midrule
		\texttt{linear} & $a^\top b$ \\
		\texttt{poly} & $(\gamma\, a^\top b + c_0)^{d}$ \\
		\texttt{rbf} / \texttt{gaussian} & $\exp(-\gamma\lVert a-b\rVert^2)$ \\
		\texttt{sigmoid} & $\tanh(\gamma\, a^\top b + c_0)$ \\
		\bottomrule
	\end{tabular}
\end{table}

\lstinputlisting[caption={Parameter dan fungsi kernel (\texttt{src/models/svm.py})},linerange={5-5,7-10,12-23,25-28}]{../src/models/svm.py}

\subsection{Sequential Minimal Optimization}

Masalah dual diselesaikan dengan \textit{Sequential Minimal Optimization}
(SMO). Setiap iterasi SMO hanya mengubah dua variabel, yaitu
$\alpha_i$ dan $\alpha_j$. Subproblem dua variabel ini dapat diselesaikan secara
analitik tanpa melanggar $\sum_i\alpha_i y_i=0$. Pasangan $(i,j)$ dipilih dengan
\textit{second-order working set selection} melalui langkah
berikut.
\begin{enumerate}
	\item Untuk setiap sampel, simpan $u_k=\sum_l\alpha_l y_l K(x_l,x_k)$ dan
	$g_k=y_k-u_k$. Himpunan $I_{\text{up}}$ berisi sampel yang $\alpha$-nya masih
	dapat bergerak searah $y_k$ positif. Himpunan $I_{\text{low}}$ berisi sampel
	yang $\alpha$-nya dapat bergerak ke arah sebaliknya.
	\item Pilih $i=\arg\max_{k\in I_{\text{up}}} g_k$, yaitu sampel yang paling
	melanggar kondisi KKT.
	\item Hentikan pelatihan jika $g_i-\min_{k\in I_{\text{low}}}g_k<\texttt{tol}$
	karena kondisi KKT sudah terpenuhi.
	\item Di antara sampel $I_{\text{low}}$ dengan $g_j<g_i$, pilih $j$ yang
	memaksimalkan kenaikan objektif $(g_i-g_j)^2/(K_{ii}+K_{jj}-2K_{ij})$.
	\item Perbarui $\alpha_j$ dengan langkah Newton
	$\alpha_j+y_j(g_j-g_i)/\eta$, dengan $\eta=K_{ii}+K_{jj}-2K_{ij}$, lalu potong
	hasilnya ke batas $[L,H]$ agar $0\le\alpha\le C$. Sesuaikan $\alpha_i$ agar
	$\sum\alpha y$ tidak berubah.
	\item Perbarui $u$ secara inkremental dengan baris kernel ke-$i$ dan ke-$j$.
\end{enumerate}
Setiap iterasi hanya menghitung dua baris kernel, sedangkan matriks kernel penuh
berukuran $n\times n$ tidak pernah disimpan. Untuk 21.857 baris data latih, matriks
penuh dalam \textit{float64} membutuhkan sekitar 3,8 GB memori. Untuk kernel
Gaussian, norma kuadrat setiap baris (\texttt{sq}) dihitung sekali di awal
sehingga satu baris kernel cukup dihitung dengan satu perkalian matriks-vektor.
Fungsi \texttt{snap} membulatkan nilai $\alpha$ yang sangat dekat dengan 0 atau
$C$ agar status \textit{bound} tidak berganti-ganti akibat galat
\textit{floating point}.

Setelah konvergen, bias $b$ dihitung sebagai rata-rata $g_k$ pada \textit{free
support vector} ($0<\alpha_k<C$) karena sampel tersebut terletak tepat di margin.
Untuk kernel linear, bobot $w=\sum_i\alpha_i y_i x_i$ disimpan di
\texttt{coef\_}. Prediksi kernel linear cukup dihitung sebagai $Xw+b$, dan nilai
$|w|$ dipakai untuk seleksi fitur (Subbab~\ref{sec:seleksi}).

\lstinputlisting[caption={Pelatihan SVM dengan SMO (\texttt{src/models/svm.py})},linerange={30-72}]{../src/models/svm.py}

\lstinputlisting[caption={Fungsi keputusan dan prediksi SVM (\texttt{src/models/svm.py})},linerange={75-83}]{../src/models/svm.py}

\section{Decision Tree Learning}
\label{sec:dtl}

\todo{Implementasi \texttt{DTL.fit} dan \texttt{DTL.predict} di
\texttt{src/models/dtl.py} belum selesai. Parameter yang sudah disiapkan:
\texttt{criterion} (\texttt{entropy}/\texttt{gini}), \texttt{max\_depth}, dan
\texttt{min\_samples\_split}. Jelaskan perhitungan \textit{information gain},
penanganan atribut kontinu, kriteria berhenti, dan potongan kode. Jika DTL
menyediakan \texttt{feature\_importances\_}, \texttt{Pipeline} otomatis dapat
memakainya untuk seleksi fitur.}

\section{Bayesian Network}
\label{sec:bn}

\todo{Implementasi \texttt{BayesNetwork.fit} dan \texttt{predict} di
\texttt{src/models/bayes\_network.py} belum selesai. Parameter yang sudah
disiapkan: \texttt{n\_bins} untuk diskretisasi fitur kontinu. Jelaskan cara
membangun struktur jaringan, estimasi \textit{conditional probability table}
(beserta \textit{smoothing}), dan inferensi posterior.}

\section{Validasi dan Utilitas Eksperimen}

Seluruh evaluasi dijalankan melalui fungsi di \texttt{src/validation.py} agar
semua model diuji dengan prosedur yang sama. Fungsi yang tersedia adalah sebagai
berikut.
\begin{itemize}
	\item \texttt{holdout} melatih model pada \texttt{X\_train} dan mengevaluasinya
	pada \texttt{X\_val}.
	\item \texttt{kfold} menjalankan \textit{stratified 5-fold cross validation}
	pada \texttt{X\_train} serta melaporkan rata-rata dan simpangan baku
	\textit{macro F1}.
	\item \texttt{tune} menjalankan \textit{grid search} dengan skor
	\textit{k-fold}, lalu \texttt{best} mengambil parameter terbaik.
	\item \texttt{resample} menerapkan SMOTE atau
	\textit{random undersampling} hanya pada data latih di setiap \textit{fold}.
	Data validasi tidak pernah berisi sampel sintetis.
	\item \texttt{confusion}, \texttt{report}, \texttt{misclassified}, dan
	\texttt{importance} dipakai untuk analisis kesalahan.
	\item \texttt{submit} melatih ulang model pada seluruh data berlabel, menulis
	prediksi ke \texttt{output/<model>\_<timestamp>.csv}, dan menyimpan model ke
	berkas \texttt{.pkl}.
\end{itemize}
Setiap model di-\texttt{deepcopy} sebelum dilatih sehingga objek model di
\textit{notebook} tidak berubah setelah validasi. Model yang belum
diimplementasikan (\texttt{NotImplementedError}) dilewati secara otomatis.

\lstinputlisting[caption={\textit{Hold-out} dan \textit{k-fold} (\texttt{src/validation.py})},linerange={13-58}]{../src/validation.py}
